% Bai 03
\chapter{Bài 3: Kỹ thuật lập trình Python (NumPy, Matplotlib)}

% Gợi ý: bạn thay thế/điền các bài đúng theo notebook `python_advance.ipynb` hoặc `my_lesson3.ipynb`.

\section{Bài 3.1: Tính toán trên dãy số với NumPy}
\subsection{Mô tả bài tập và ý tưởng giải quyết}
\begin{itemize}
  \item Tạo mảng NumPy \(a\) gồm \(N=10\) phần tử.
  \item Thực hiện các phép toán: bình phương, tìm max/min/mean, tính phương sai và độ lệch chuẩn.
  \item In kết quả ra màn hình để kiểm tra.
\end{itemize}

\subsection{Code cài đặt}
\begin{lstlisting}
import numpy as np

a = np.array([120, 183, 160, 175, 145, 162, 190, 160, 152, 162])
a2 = a * a

v_max = np.max(a)
v_min = np.min(a)
v_avg = np.average(a)

v_mean = np.sum(a) / len(a)
v_sigma2 = np.sum((a - v_mean) ** 2) / (len(a) - 1)
v_var = np.sqrt(v_sigma2)

print("Day so a:", a)
print("Binh phuong a:", a2)
print("Max:", v_max, "Min:", v_min, "Mean:", v_avg)
print("Phuong sai:", v_sigma2, "Do lech chuan:", v_var)
\end{lstlisting}

\subsection{Kết quả chạy và giải thích}
\ScreenshotFigure{assets/screenshots/bai03_b1_numpy_array_ops.png}{Kết quả thao tác mảng NumPy (screenshot).}

\textbf{Giải thích:} NumPy hỗ trợ thao tác vector hoá nên các phép tính max/min/mean/variance thực hiện nhanh và gọn.

\section{Bài 3.2: Vẽ đồ thị với Matplotlib}
\subsection{Mô tả bài tập và ý tưởng giải quyết}
\begin{itemize}
  \item Sinh dữ liệu hàm số (ví dụ: \(\sin(x)\), \(\cos(x)\)) trong khoảng \(x\in[-\pi,\pi]\).
  \item Vẽ 2 đường trên cùng biểu đồ, gắn nhãn và hiển thị.
  \item Chụp lại hình kết quả từ notebook.
\end{itemize}

\subsection{Code cài đặt}
\begin{lstlisting}
import numpy as np
import matplotlib.pyplot as plt

X = np.linspace(-np.pi, np.pi, 256, endpoint=True)
C, S = np.cos(X), np.sin(X)

plt.plot(X, C, label="cos(x)")
plt.plot(X, S, label="sin(x)")
plt.legend()
plt.grid(True)
plt.show()
\end{lstlisting}

\subsection{Kết quả chạy và giải thích}
\ScreenshotFigure{assets/screenshots/bai03_b2_plot_sin_cos.png}{Đồ thị \(\sin(x)\) và \(\cos(x)\) (screenshot).}

\textbf{Giải thích:} Đồ thị thể hiện đúng tính chất tuần hoàn của hai hàm lượng giác và sự lệch pha giữa \(\sin\) và \(\cos\).

\clearpage

